% Created 2026-10-05 Mon 16:30
% Intended LaTeX compiler: xelatex
\documentclass[12pt]{article}
\usepackage{graphicx}
\usepackage{longtable}
\usepackage{wrapfig}
\usepackage{rotating}
\usepackage[normalem]{ulem}
\usepackage{amsmath}
\usepackage{amssymb}
\usepackage{capt-of}
\usepackage{hyperref}
\usepackage[a4paper,margin=1in]{geometry}
\usepackage{xurl}
\usepackage{graphicx}
\usepackage{parskip}
\usepackage{titling}
\makeatletter
\renewcommand{\maketitle}{%
\begin{center}
{\LARGE\@title\par}
\vspace{0.5em}
{\large\@author\par}
\vspace{0.5em}
{\large\@date\par}
\end{center}
}
\makeatother
\author{Filipa Calado}
\date{}
\title{Service Statement}
\hypersetup{
 pdfauthor={Filipa Calado},
 pdftitle={Service Statement},
 pdfkeywords={},
 pdfsubject={},
 pdfcreator={Emacs 29.3 (Org mode 9.6.15)}, 
 pdflang={English}}
\begin{document}

\maketitle
This service statement covers my service work at the Pratt Institute, School of
Information (SI), at Pratt Institute broadly, and my external service. 

\section*{Pratt SI}
\label{sec:org626d777}
Since beginning my position at Pratt School of Information (SI), I
have been involved in faculty committees and taskforces, student
groups, and in the faculty union. I have also been a faculty advisor
to 30 - 35 students per semester, and ran workshops for the Digital
Storytelling Lab in AY 2025-2026 and 2026-2027.

From my faculty committees and curriculum development, I want to
highlight my work in the DH Certificate Review (in collaboration with
my colleague, John Decker), the Ethics and Technology Forum Committee,
the MS Critical and Creative Computing Task Force, the DEI Committee,
and my participation on two faculty searches, the Pratt Faculty
Council, and the Peer Review Committee. Please see below for a more
exhaustive list of my service contributions.

I also want to highlight my work as a Faculty Advisor, in particular
to the new student group, the \emph{Stacks} Student Journal. This group
emerged from my INFO 601 course, and they currently are working toward
their first issue of the journal that will showcase the high level of
intellectual work being accomplished in the 601 course.

Below is an exhaustive list of my service at Pratt SI:
\begin{itemize}
\item Responsible AI Task Force (2026-2027)
\item MS Critical and Creative Computing Task Force (2026-2027)
\item Ethics and Technology Forum Committee (Summer 2026)
\item DH Certification Review (2025-2027)
\item DEI Committee (2025-2027)
\item Faculty Search Committee (2024-2026)
\item Peer Review Committee (2024-2027)
\item Faculty Council (2024-2027)
\item Faculty Advisor, \emph{Stacks} Student Journal (2025-2027)
\item Faculty Advisor, Latines in Tech (2025-2027)
\item Speaker, Digital Spring Clean-Up Day event (Spring 2026)
\end{itemize}

\section*{Pratt Institute}
\label{sec:org58a18f7}
Beyond SI, I have also been involved in the Pratt Faculty Union,
UFCT 1460.

With my colleague Monica Maceli, I am currently a Union Delegate for
the School of Information. Besides that, I worked on the Summer 2026
membership drive, an initiative to increase membership at the union
which involved one-on-one organizing conversations with several of my
colleagues at Pratt SI.

This summer, I joined the Data Team with Philip Ording, my colleague
from the Math and Science department at Pratt. Together, Philip and I
are working to establish a new database and data workflows for union
data, with the goal of preparing our membership for the next round of
bargaining, in 2027.

A full list of my tasks and roles for the union includes the
following:
\begin{itemize}
\item Member of the Data Team (2026-2027)
\item Union Delegate for School of Information (2025-2027)
\item Participant in the Summer 2026 membership drive (Summer 2026)
\item Participant in ``Organizing For Power'' training series
(Spring 2025)
\end{itemize}

\section*{External}
\label{sec:org7181e17}
My goal with external service obligations is to maintain individual
committments to humanities and computer science scholarly communities.
For that purpose, I serve as a reviewer for two publications.

The first publication is \emph{\href{https://read.dukeupress.edu/social-text}{\uline{Social Text}}}, published by Duke University
Press. Known for its radical and activist leaning, \emph{Social Text} is a
leading peer-reviewed journal in cultural studies, one that has
published consistently provocative work on gender, sexuality, race,
colonialism, environmental issues, and popular culture for nearly 50
years. I have been a member of the journal's editorial board for three
years, regularly attending their quarterly meetings and reviewing
submissions several times per year. In addition to participating in
the collective, I have also joined my colleagues Jayna Browne (Brown
University), and Ángeles Donoso Macaya (CUNY Graduate Center) to form
the Denkbild review committee, for reviewing short-form submissions to
the journal.

The second publication is the \uline{\href{https://gebnlp-workshop.github.io/}{Workshop on Gender Bias in Natural
Language Processing}} (GeBNLP), held annually at different conference
venues. GeBNLP is an annual workshop that takes place every year for
the past 6 years. At last year's conference, held at the Association
for Computational Linguistics (ACL) in Vienna, I presented my paper,
"Some Myths About Bias: A Queer Studies Reading Of Gender Bias In NLP"
in their \emph{Proceedings of the 6th Workshop on Gender Bias in Natural
Language Processing (GeBNLP)}. For this year's workshop, taking place
at AACL-IJCNLP 2026, I am participating as a reviewer. I plan to
continue as both a reviewer and contributer to the workshop in future
years.
\end{document}
